% RQ3: Per-shot policy trajectories and backlog-cost comparison for two devices.
%
% S26 Ultra currently renders the collected No-pacing and Ours series; empty
% policy CSVs add no curve until their factual traces arrive.  S26 remains Data
% pending.
%
% Each policy CSV contains one row per shot and the median, P10, and P90 of
% realized backlog, queue depth, and applied delay.  Each backlog-cost CSV may
% contain multiple parameter settings per policy to draw a frontier.

\begin{figure*}[t]
  \centering
  \newif\ifrqthreeultrahasdata
  \newif\ifrqthreesixhasdata
  \providecommand{\rqThreeUltraDataSwitch}{\rqthreeultrahasdatatrue}
  \providecommand{\rqThreeStandardDataSwitch}{\rqthreesixhasdatafalse}
  \providecommand{\rqThreeDataDir}{data/rq3}
  \providecommand{\rqThreeUltraDataDir}{\rqThreeDataDir/s26_ultra}
  \providecommand{\rqThreeStandardDataDir}{\rqThreeDataDir/s26}
  \rqThreeUltraDataSwitch
  \rqThreeStandardDataSwitch

  % Apply the same explicit bounds to both device rows after data collection.
  \providecommand{\rqThreeBacklogAxis}{ymin=0}
  \providecommand{\rqThreeQueueAxis}{ymin=0}
  \providecommand{\rqThreeDelayAxis}{ymin=0}
  \providecommand{\rqThreeFrontierAxis}{xmin=0,ymin=0}

  \tikzset{
    rq3 no pacing/.style={gray, densely dotted, line width=0.8pt},
    rq3 static/.style={blue!70!black, dashed, line width=0.8pt},
    rq3 dynamic/.style={red!75!black, dash dot, line width=0.8pt},
    rq3 ours/.style={black, solid, line width=1.15pt},
  }

  \newcommand{\rqthreepending}{%
    \node[font=\scriptsize,text=gray] at (rel axis cs:0.5,0.5)
      {Data pending};}

  % #1: device data directory; #2: metric column prefix.
  \newcommand{\rqthreepolicycurves}[2]{%
    \addplot[rq3 no pacing,forget plot]
      table[x=shot,y=#2_median,col sep=comma]
      {#1/no_pacing.csv};
    \addplot[rq3 static,forget plot]
      table[x=shot,y=#2_median,col sep=comma]
      {#1/static_lut.csv};
    \addplot[rq3 dynamic,forget plot]
      table[x=shot,y=#2_median,col sep=comma]
      {#1/queue_dynamic.csv};
    \addplot[rq3 ours,forget plot]
      table[x=shot,y=#2_median,col sep=comma]
      {#1/ours.csv};
  }

  % #1: device data directory.
  \newcommand{\rqthreefrontiercurves}[1]{%
    \addplot[rq3 no pacing,mark=o]
      table[
        x=no_pacing_delay_s,
        y=no_pacing_max_backlog_s,
        col sep=comma
      ] {#1/backlog_cost.csv};
    \addplot[rq3 static,mark=square]
      table[
        x=static_delay_s,
        y=static_max_backlog_s,
        col sep=comma
      ] {#1/backlog_cost.csv};
    \addplot[rq3 dynamic,mark=triangle]
      table[
        x=queue_dynamic_delay_s,
        y=queue_dynamic_max_backlog_s,
        col sep=comma
      ] {#1/backlog_cost.csv};
    \addplot[rq3 ours,mark=*]
      table[
        x=ours_delay_s,
        y=ours_max_backlog_s,
        col sep=comma
      ] {#1/backlog_cost.csv};}

  % Keep conditionals outside pgfplots axis bodies: \addplot scans its input,
  % so a conditional spanning a plot command can be both slow and fragile.
  \ifrqthreeultrahasdata
    \newcommand{\rqthreeultrapendingaxis}{}
    \newcommand{\rqthreeultrapendingfrontieraxis}{}
    \newcommand{\rqthreeultrapaneldata}[1]{%
      \rqthreepolicycurves{\rqThreeUltraDataDir}{#1}}
    \newcommand{\rqthreeultrafrontierdata}{%
      \rqthreefrontiercurves{\rqThreeUltraDataDir}}
  \else
    \newcommand{\rqthreeultrapendingaxis}{%
      ymin=0,ymax=1,yticklabels=\empty}
    \newcommand{\rqthreeultrapendingfrontieraxis}{%
      xmax=1,ymin=0,ymax=1,xticklabels=\empty,yticklabels=\empty}
    \newcommand{\rqthreeultrapaneldata}[1]{\rqthreepending}
    \newcommand{\rqthreeultrafrontierdata}{\rqthreepending}
  \fi

  \ifrqthreesixhasdata
    \newcommand{\rqthreesixpendingaxis}{}
    \newcommand{\rqthreesixpendingfrontieraxis}{}
    \newcommand{\rqthreesixpaneldata}[1]{%
      \rqthreepolicycurves{\rqThreeStandardDataDir}{#1}}
    \newcommand{\rqthreesixfrontierdata}{%
      \rqthreefrontiercurves{\rqThreeStandardDataDir}}
  \else
    \newcommand{\rqthreesixpendingaxis}{%
      ymin=0,ymax=1,yticklabels=\empty}
    \newcommand{\rqthreesixpendingfrontieraxis}{%
      xmax=1,ymin=0,ymax=1,xticklabels=\empty,yticklabels=\empty}
    \newcommand{\rqthreesixpaneldata}[1]{\rqthreepending}
    \newcommand{\rqthreesixfrontierdata}{\rqthreepending}
  \fi

  \newsavebox{\rqthreefigurebox}
  \begin{lrbox}{\rqthreefigurebox}
  \begin{tikzpicture}
    \begin{groupplot}[
      group style={
        group size=4 by 2,
        horizontal sep=0.024\textwidth,
        vertical sep=0.060\textwidth,
      },
      width=0.215\textwidth,
      height=0.190\textwidth,
      xmin=1,
      xmax=30,
      xtick={1,10,20,30},
      tick label style={font=\tiny},
      label style={font=\scriptsize},
      title style={font=\scriptsize},
      grid=major,
      major grid style={gray!20,line width=0.2pt},
      axis line style={line width=0.35pt},
      tick style={line width=0.35pt},
      unbounded coords=discard,
      enlarge y limits=0.08,
    ]

      \nextgroupplot[
        title={(a) Backlog (ms)},
        xticklabels=\empty,
        \rqThreeBacklogAxis,
        \rqthreeultrapendingaxis
      ]
      \rqthreeultrapaneldata{backlog}

      \nextgroupplot[
        title={(b) Queue depth},
        xticklabels=\empty,
        \rqThreeQueueAxis,
        \rqthreeultrapendingaxis
      ]
      \rqthreeultrapaneldata{queue_depth}

      \nextgroupplot[
        title={(c) Pacing delay (ms)},
        xticklabels=\empty,
        \rqThreeDelayAxis,
        \rqthreeultrapendingaxis
      ]
      \rqthreeultrapaneldata{delay}

      \nextgroupplot[
        title={(d) Backlog vs.\ delay},
        \rqThreeFrontierAxis,
        \rqthreeultrapendingfrontieraxis
      ]
      \rqthreeultrafrontierdata

      \nextgroupplot[
        title={(e) Backlog (ms)},
        xlabel={Shot index},
        \rqThreeBacklogAxis,
        \rqthreesixpendingaxis
      ]
      \rqthreesixpaneldata{backlog}

      \nextgroupplot[
        title={(f) Queue depth},
        xlabel={Shot index},
        \rqThreeQueueAxis,
        \rqthreesixpendingaxis
      ]
      \rqthreesixpaneldata{queue_depth}

      \nextgroupplot[
        title={(g) Pacing delay (ms)},
        xlabel={Shot index},
        \rqThreeDelayAxis,
        \rqthreesixpendingaxis
      ]
      \rqthreesixpaneldata{delay}

      \nextgroupplot[
        title={(h) Backlog vs.\ delay},
        xlabel={Cumulative delay (s)},
        \rqThreeFrontierAxis,
        \rqthreesixpendingfrontieraxis
      ]
      \rqthreesixfrontierdata
    \end{groupplot}

    \node[font=\scriptsize\bfseries,anchor=east,align=center]
      at ([xshift=-1.2ex]group c1r1.west) {S26\\Ultra};
    \node[font=\scriptsize\bfseries,anchor=east]
      at ([xshift=-1.2ex]group c1r2.west) {S26};

    % One shared policy legend applies to both device rows.
    \coordinate (rqthreelegendpos) at
      ($(group c1r2.south)!0.5!(group c4r2.south)$);
    \matrix[
      matrix of nodes,
      nodes={font=\scriptsize,anchor=west},
      column sep=2pt,
      row sep=0pt,
      anchor=north,
    ] at ([yshift=-4.0ex]rqthreelegendpos) {
      \draw[rq3 no pacing] (0,0)--(0.38,0); & No pacing &
      \draw[rq3 static] (0,0)--(0.38,0); & Static LUT &
      \draw[rq3 dynamic] (0,0)--(0.38,0); & Queue dynamic &
      \draw[rq3 ours] (0,0)--(0.38,0); & Ours \\
    };
  \end{tikzpicture}
  \end{lrbox}
  \resizebox{\textwidth}{!}{\usebox{\rqthreefigurebox}}

  \caption{RQ3 pacing behavior on S26 Ultra and S26 under 12MP normal at
  starting thermal level 3.  Rows separate devices; columns report per-shot
  median backlog, queue depth, pacing delay, and the backlog--delay comparison.
  Available factual policies are plotted; uncollected policies remain absent.
  The current S26 Ultra No-pacing and Ours curves are provisional because their
  admitted workloads were not matched.}
  \label{fig:rq3_pacing_trajectories}
\end{figure*}
